\documentclass[11pt]{article}
\usepackage{amsmath,amssymb,amsthm,mathtools}
\usepackage[a4paper,margin=1in]{geometry}
\usepackage{enumitem}
\usepackage{hyperref}
\usepackage{framed}

\title{Lecture Notes: Concrete Cohomology on \(S^2\) and \(T^2\)\\
from Winding Forms to Gauss--Bonnet}
\author{}
\date{}

\newtheorem{theorem}{Theorem}
\newtheorem{remark}{Remark}
\newtheorem{exercise}{Exercise}

\newenvironment{takeaway}
{\begin{framed}\noindent\textbf{Takeaway. }}
{\end{framed}}

\newenvironment{remember}
{\begin{framed}\noindent\textbf{What should I remember? }}
{\end{framed}}

\begin{document}
\maketitle

\tableofcontents

\section{The main question}

These notes are about one question:

\[
\boxed{\text{When is a closed form not exact?}}
\]

That is the concrete meaning of cohomology.

We will study this through three examples:
\begin{itemize}[leftmargin=2em]
\item the winding form on \(\mathbb{R}^2\setminus\{0\}\),
\item the area form on \(S^2\),
\item the forms \(dx,dy,dx\wedge dy\) on \(T^2\).
\end{itemize}

Then we will connect this to curvature and Gauss--Bonnet.

\begin{takeaway}
Cohomology measures the failure of local primitives to glue into global primitives.
\end{takeaway}

\section{de Rham cohomology}

For a smooth manifold \(M\),
\[
H^k_{\mathrm{dR}}(M)
=
\frac{\{\text{closed }k\text{-forms}\}}{\{\text{exact }k\text{-forms}\}}.
\]

So:
\begin{itemize}[leftmargin=2em]
\item \(\omega\) is \emph{closed} if \(d\omega=0\),
\item \(\omega\) is \emph{exact} if \(\omega=d\eta\).
\end{itemize}

Every exact form is closed:
\[
\omega=d\eta \implies d\omega=d^2\eta=0.
\]

But not every closed form is exact.

That is exactly where cohomology lives.

\begin{takeaway}
A nonzero cohomology class is represented by a form that is closed but not exact.
\end{takeaway}

\begin{remember}
To understand cohomology concretely, do not start with abstraction. Start with:
\[
\text{compute } d\omega,\quad \text{compute periods/integrals},\quad \text{test exactness.}
\]
\end{remember}

\section{First example: the winding form on \(\mathbb{R}^2\setminus\{0\}\)}

Consider
\[
\eta=\frac{-y\,dx+x\,dy}{x^2+y^2}.
\]

This is the standard winding form.

\subsection{Step 1: rewrite in polar coordinates}

Write
\[
x=r\cos\theta,\qquad y=r\sin\theta.
\]
Then
\[
dx=\cos\theta\,dr-r\sin\theta\,d\theta,
\qquad
dy=\sin\theta\,dr+r\cos\theta\,d\theta.
\]

Compute:
\begin{align*}
-y\,dx+x\,dy
&=
-r\sin\theta(\cos\theta\,dr-r\sin\theta\,d\theta)
+r\cos\theta(\sin\theta\,dr+r\cos\theta\,d\theta)\\
&=
-r\sin\theta\cos\theta\,dr+r^2\sin^2\theta\,d\theta
+r\sin\theta\cos\theta\,dr+r^2\cos^2\theta\,d\theta\\
&=
r^2(\sin^2\theta+\cos^2\theta)\,d\theta\\
&=
r^2\,d\theta.
\end{align*}

Therefore
\[
\eta=d\theta.
\]

So locally the winding form is just the differential of the angle.

\subsection{Step 2: why it is closed}

Since locally \(\eta=d\theta\),
\[
d\eta=d(d\theta)=0.
\]

So \(\eta\) is closed.

\subsection{Step 3: why it is not exact}

Integrate around the unit circle:
\[
\gamma(t)=(\cos t,\sin t),\qquad 0\le t\le 2\pi.
\]
Then
\[
\int_\gamma \eta=\int_0^{2\pi} d\theta=2\pi.
\]

But if \(\eta=df\) globally, then for every closed curve,
\[
\int_\gamma \eta=\int_\gamma df=0.
\]

Contradiction.

So \(\eta\) is not exact.

Thus
\[
[\eta]\neq 0 \in H^1_{\mathrm{dR}}(\mathbb{R}^2\setminus\{0\}).
\]

\begin{takeaway}
The winding form is closed but not exact because there is no globally single-valued angle function on \(\mathbb{R}^2\setminus\{0\}\).
\end{takeaway}

\begin{remember}
The first concrete test for non-exactness is:
\[
\text{find a closed curve with nonzero integral.}
\]
\end{remember}

\section{Second example: the sphere \(S^2\)}

Now move from a \(1\)-form on a punctured plane to a \(2\)-form on a sphere.

The key object is the area form.

\subsection{Step 1: spherical coordinates}

Write
\[
X(\theta,\phi)=
(\sin\phi\cos\theta,\sin\phi\sin\theta,\cos\phi),
\qquad
0<\theta<2\pi,\quad 0<\phi<\pi.
\]

The induced metric is
\[
ds^2=d\phi^2+\sin^2\phi\,d\theta^2.
\]

So the area form is
\[
dA=\sin\phi\,d\theta\wedge d\phi.
\]

\subsection{Step 2: why \(dA\) is closed}

Since \(S^2\) is \(2\)-dimensional, any \(3\)-form on \(S^2\) is automatically zero.

But \(d(dA)\) is a \(3\)-form.

So
\[
d(dA)=0.
\]

Therefore \(dA\) is closed.

\subsection{Step 3: compute the total integral}

Now compute:
\begin{align*}
\int_{S^2}dA
&=
\int_0^{2\pi}\int_0^\pi \sin\phi\,d\phi\,d\theta\\
&=
\int_0^{2\pi}\left[-\cos\phi\right]_{\phi=0}^{\phi=\pi}\,d\theta\\
&=
\int_0^{2\pi}2\,d\theta\\
&=
4\pi.
\end{align*}

\subsection{Step 4: why \(dA\) is not exact}

If \(dA=d\alpha\) globally, then by Stokes,
\[
\int_{S^2}dA=\int_{S^2}d\alpha=\int_{\partial S^2}\alpha=0.
\]

But we just found
\[
\int_{S^2}dA=4\pi.
\]

So \(dA\) is not exact.

Thus
\[
[dA]\neq 0\in H^2_{\mathrm{dR}}(S^2).
\]

\begin{takeaway}
On the sphere, the interesting nontrivial cohomology class is in degree \(2\), represented by the area form.
\end{takeaway}

\begin{remember}
For a compact oriented surface,
\[
\text{if a }2\text{-form has nonzero total integral, it cannot be exact.}
\]
\end{remember}

\section{Why \(H^1(S^2)=0\)}

This is worth pausing over.

On \(\mathbb{R}^2\setminus\{0\}\), the hole gave a nontrivial \(H^1\)-class.

On \(S^2\), there is no such \(1\)-dimensional hole.

Concretely, every loop on \(S^2\) bounds a region. So if \(\omega\) is a closed \(1\)-form and \(\gamma=\partial D\), then
\[
\int_\gamma \omega = \int_D d\omega = 0.
\]

So all periods vanish.

Hence every closed \(1\)-form is exact, and
\[
H^1_{\mathrm{dR}}(S^2)=0.
\]

Therefore
\[
H^0_{\mathrm{dR}}(S^2)\cong\mathbb{R},\qquad
H^1_{\mathrm{dR}}(S^2)=0,\qquad
H^2_{\mathrm{dR}}(S^2)\cong\mathbb{R}.
\]

\begin{takeaway}
The sphere has no winding \(1\)-classes, but it does have a nontrivial top-degree area class.
\end{takeaway}

\section{Local primitives on \(S^2\)}

This is the sphere-version of the winding story.

Cover \(S^2\) by
\[
U_N=S^2\setminus\{N\},\qquad
U_S=S^2\setminus\{S\}.
\]

Define
\[
\alpha_N=(1-\cos\phi)\,d\theta
\qquad\text{on }U_N,
\]
and
\[
\alpha_S=-(1+\cos\phi)\,d\theta
\qquad\text{on }U_S.
\]

Then
\[
d\alpha_N=\sin\phi\,d\phi\wedge d\theta,
\qquad
d\alpha_S=\sin\phi\,d\phi\wedge d\theta.
\]

So on both charts, the area form is locally a derivative.

Now subtract on the overlap:
\[
\alpha_N-\alpha_S
=
(1-\cos\phi)d\theta+(1+\cos\phi)d\theta
=
2\,d\theta.
\]

So the mismatch between the two local primitives is twice the winding form.

\begin{takeaway}
The area form on \(S^2\) is locally exact, but the local primitives differ by angular winding on the overlap.
\end{takeaway}

\begin{remember}
For \(S^2\):
\[
\text{local exactness} + \text{overlap winding} = \text{global non-exactness}.
\]
\end{remember}

\section{Third example: the torus \(T^2\)}

Now we study
\[
T^2=\mathbb{R}^2/\mathbb{Z}^2.
\]

Unlike the sphere, the torus really does have two independent winding directions.

\subsection{Step 1: the forms \(dx\) and \(dy\)}

The coordinate forms \(dx\) and \(dy\) on \(\mathbb{R}^2\) descend to global \(1\)-forms on \(T^2\).

They are closed:
\[
d(dx)=0,\qquad d(dy)=0.
\]

\subsection{Step 2: test exactness by periods}

Consider the basic loops
\[
\gamma_x(t)=(t,0)\mod\mathbb{Z}^2,\qquad
\gamma_y(t)=(0,t)\mod\mathbb{Z}^2,\qquad 0\le t\le 1.
\]

Then
\[
\int_{\gamma_x}dx=1,\qquad \int_{\gamma_x}dy=0,
\]
\[
\int_{\gamma_y}dx=0,\qquad \int_{\gamma_y}dy=1.
\]

Therefore neither \(dx\) nor \(dy\) is exact.

So
\[
[dx],[dy]\neq 0 \in H^1_{\mathrm{dR}}(T^2).
\]

And since the two loops separate them, they are linearly independent.

Hence
\[
H^1_{\mathrm{dR}}(T^2)\cong \mathbb{R}^2,
\]
with basis \([dx],[dy]\).

\begin{takeaway}
The torus has two independent winding \(1\)-classes: \([dx]\) and \([dy]\).
\end{takeaway}

\subsection{Step 3: the top class \(dx\wedge dy\)}

Now consider
\[
dx\wedge dy.
\]

It is closed:
\[
d(dx\wedge dy)=0.
\]

And
\[
\int_{T^2}dx\wedge dy
=
\int_0^1\int_0^1 dx\,dy
=
1.
\]

So it is not exact.

Thus
\[
[dx\wedge dy]\neq 0 \in H^2_{\mathrm{dR}}(T^2).
\]

So the torus has
\[
H^0_{\mathrm{dR}}(T^2)\cong\mathbb{R},\qquad
H^1_{\mathrm{dR}}(T^2)\cong\mathbb{R}^2,\qquad
H^2_{\mathrm{dR}}(T^2)\cong\mathbb{R}.
\]

\begin{takeaway}
On \(T^2\), the top-degree class comes from wedging the two basic winding classes:
\[
[dx]\wedge[dy]=[dx\wedge dy].
\]
\end{takeaway}

\begin{remember}
The torus is different from the sphere because it has genuine \(1\)-dimensional winding directions.
\end{remember}

\section{The local-to-global point of view}

Suppose \(\omega\) is a closed \(2\)-form on a manifold.

Then on each small open set \(U_i\), the Poincar\'e lemma gives
\[
\omega|_{U_i}=d\alpha_i.
\]

So closed forms always have local primitives.

On overlaps,
\[
d(\alpha_i-\alpha_j)=0.
\]

So the differences are closed \(1\)-forms, and often locally exact:
\[
\alpha_i-\alpha_j=df_{ij}.
\]

These overlap functions \(f_{ij}\) measure the failure of the local primitives to glue.

Exponentiating them gives transition functions
\[
g_{ij}=e^{if_{ij}},
\]
which define line bundles.

\begin{takeaway}
Cohomology, overlap data, and line bundles are different languages for the same obstruction.
\end{takeaway}

\section{Curvature and Gauss--Bonnet}

Now bring geometry into the picture.

Let \(M\) be an oriented Riemannian surface, with Gaussian curvature \(K\) and area form \(dA\).

Then the curvature \(2\)-form is
\[
K\,dA.
\]

\subsection{Why this is closed}

Since \(M\) is \(2\)-dimensional,
\[
d(K\,dA)=0.
\]

So \(K\,dA\) defines a cohomology class in \(H^2_{\mathrm{dR}}(M)\).

\subsection{Why this is locally exact}

Locally, if we choose a positively oriented orthonormal frame, the Levi-Civita connection is given by a local \(1\)-form \(\omega_{12}\), and
\[
d\omega_{12}=K\,dA.
\]

So curvature is locally the derivative of a connection form.

Again:
\[
\text{local primitive exists, global primitive may fail.}
\]

That failure is the Euler class.

\subsection{Gauss--Bonnet}

The Gauss--Bonnet theorem says
\[
\int_M K\,dA = 2\pi\chi(M).
\]

So total curvature is topological.

\begin{takeaway}
The cohomology class of the curvature form remembers the Euler characteristic.
\end{takeaway}

\begin{remember}
Gauss--Bonnet is not about the area form \(dA\) by itself. It is about the curvature form
\[
K\,dA.
\]
\end{remember}

\section{Check Gauss--Bonnet on \(S^2\)}

For the unit sphere,
\[
K\equiv 1.
\]

So
\[
K\,dA=dA.
\]

Hence
\[
\int_{S^2}K\,dA
=
\int_{S^2}dA
=
4\pi.
\]

Since
\[
\chi(S^2)=2,
\]
we get
\[
2\pi\chi(S^2)=4\pi.
\]

Therefore
\[
\int_{S^2}K\,dA=2\pi\chi(S^2).
\]

\begin{takeaway}
On \(S^2\), the curvature class is nontrivial:
\[
[K\,dA]=[dA]\neq 0.
\]
\end{takeaway}

\section{Check Gauss--Bonnet on \(T^2\)}

Take the flat torus with metric
\[
ds^2=dx^2+dy^2.
\]

Then
\[
K\equiv 0.
\]

So
\[
K\,dA=0.
\]

Hence
\[
\int_{T^2}K\,dA=0.
\]

Since
\[
\chi(T^2)=0,
\]
we get
\[
\int_{T^2}K\,dA=2\pi\chi(T^2)=0.
\]

\begin{takeaway}
On the flat torus, the curvature class vanishes even though the area class does not.
\end{takeaway}

\begin{remember}
Do not confuse:
\[
[dA]\neq 0
\qquad\text{with}\qquad
[K\,dA].
\]
For a flat torus,
\[
[dA]\neq 0,\qquad [K\,dA]=0.
\]
\end{remember}

\section{Euler class on \(S^2\) from stereographic coordinates}

We now make the sphere story even more concrete.

Let
\[
N=(0,0,1),\qquad S=(0,0,-1).
\]

Define stereographic coordinates:
\[
z=\frac{X+iY}{1-Z}
\qquad\text{from the north pole,}
\]
\[
w=\frac{X-iY}{1+Z}
\qquad\text{from the south pole.}
\]

On the overlap,
\[
w=\frac{1}{z}.
\]

So the tangent-bundle transition function is the derivative:
\[
\frac{dw}{dz}=-\frac{1}{z^2}.
\]

Restrict to the equator, where \(z=e^{i\theta}\). Then
\[
\frac{dw}{dz}=-e^{-2i\theta}.
\]

This winds twice around the unit circle as \(\theta\) runs from \(0\) to \(2\pi\).

So the tangent bundle of \(S^2\) has clutching degree \(2\), hence Euler number
\[
e(TS^2)[S^2]=2.
\]

This matches
\[
\chi(S^2)=2
\]
and also matches
\[
\frac{1}{2\pi}\int_{S^2}K\,dA
=
\frac{1}{2\pi}\cdot 4\pi
=
2.
\]

\begin{takeaway}
For the sphere, overlap winding, Euler class, and total curvature all encode the same number:
\[
2.
\]
\end{takeaway}

\section{Big picture}

Let us compare the three spaces.

\subsection*{\(\mathbb{R}^2\setminus\{0\}\)}
The winding form is closed but not exact:
\[
\eta=\frac{-y\,dx+x\,dy}{x^2+y^2},\qquad [\eta]\neq 0\in H^1.
\]

\subsection*{\(S^2\)}
The area form is closed but not exact:
\[
dA=\sin\phi\,d\theta\wedge d\phi,\qquad [dA]\neq 0\in H^2.
\]
Also,
\[
\int_{S^2}K\,dA=4\pi=2\pi\chi(S^2).
\]

\subsection*{\(T^2\)}
The classes \([dx]\) and \([dy]\) generate \(H^1\), and \([dx\wedge dy]\) generates \(H^2\). For the flat torus,
\[
[K\,dA]=0.
\]

\begin{takeaway}
The same idea appears everywhere:
\[
\boxed{\text{local primitives exist, but global primitives may fail to exist.}}
\]
\end{takeaway}

\section{Exercises}

\begin{exercise}
Show directly that
\[
d\!\left(\frac{-y\,dx+x\,dy}{x^2+y^2}\right)=0
\]
on \(\mathbb{R}^2\setminus\{0\}\).
\end{exercise}

\begin{exercise}
Compute
\[
\int_{|z|=1}\frac{-y\,dx+x\,dy}{x^2+y^2}.
\]
Why does this prove the winding form is not exact?
\end{exercise}

\begin{exercise}
Verify that on the unit sphere,
\[
dA=\sin\phi\,d\theta\wedge d\phi.
\]
\end{exercise}

\begin{exercise}
Compute
\[
\int_{S^2}dA=4\pi.
\]
Why does this prove \(dA\) is not exact?
\end{exercise}

\begin{exercise}
Compute
\[
d((1-\cos\phi)d\theta),\qquad
d(-(1+\cos\phi)d\theta),
\]
and compare the two local primitives on the overlap.
\end{exercise}

\begin{exercise}
On \(T^2=\mathbb{R}^2/\mathbb{Z}^2\), compute the periods of \(dx\) and \(dy\) on the two basic loops and prove that \([dx]\) and \([dy]\) are linearly independent.
\end{exercise}

\begin{exercise}
Compute
\[
\int_{T^2}dx\wedge dy
\]
and explain why \(dx\wedge dy\) is not exact.
\end{exercise}

\begin{exercise}
Check Gauss--Bonnet directly for the unit sphere and the flat torus.
\end{exercise}

\begin{exercise}
Explain why the stereographic transition function of \(TS^2\) has winding number \(2\) along the equator.
\end{exercise}

\end{document}